\begin{lemma}
Let $1<p<\infty$, let $\Omega\subseteq\mathbb{C}$ be open, let
$S\subseteq\Omega$, let $f\colon\mathbb{C}\to\mathbb{C}$ and let
$v\colon\mathbb{C}\to\mathbb{R}$.  Assume
\begin{enumerate}
\item[(H1)] $S\cap K$ is finite for every compact $K\subseteq\Omega$;
\item[(H2)] for $z\in\Omega$, $f(z)=0$ if and only if $z\in S$;
\item[(H3)] $f$ is continuous on $\Omega$;
\item[(H4)] $f$ is $C^\infty$ on $\Omega-S$;
\item[(H5)] $f$ solves equation (12) on $\Omega$ in the sense of
  \noderef{SolvesEq12};
\item[(H6)] $v$ is continuous on $\Omega$;
\item[(H7)] $v$ is $C^\infty$ on $\Omega-S$;
\item[(H8)] $\partial v/\partial z=f$ on $\Omega-S$, with $\partial/\partial z$
  as in \noderef{DZ};
\item[(H9)] $v$ and $f$ belong to $L^p_{\mathrm{loc}}(\Omega)$ in the sense of
  \noderef{MemLpLocOn};
\item[(H10)] both first-order partial derivatives $\partial_j w$,
  $j\in\{0,1\}$, of $w=|f|^{p-2}f$ in the sense of \noderef{PartialWord}
  belong to $L^1_{\mathrm{loc}}(\Omega)$.
\end{enumerate}
Then $v$ is $p$-harmonic on $\Omega$ in the sense of \noderef{PHarmonicOn}.

This is the converse of the quoted result (C1), and the paper's note lists it
among the facts that must be proved: ``the inverse of (54) is the complex
gradient of a $p$-harmonic function on $\mathbb{C}$''.  The potential is taken
as a hypothesis in the paper's own form $\partial v/\partial z=f$ of (55);
constructing a potential is a separate matter, which in Section 7 uses the
homogeneity of the extremal inverse.
\end{lemma}
\begin{proof}
Write $c'=\tfrac1p-\tfrac12=\tfrac{2-p}{2p}$, and let
$\varphi_v(z)$ denote $v(z)$ regarded as a complex number, so that (H8) reads
$\partial_z\varphi_v=f$ on $\Omega-S$.  Set
\[
  w(z)=|f(z)|^{p-2}f(z)\qquad(z\in\mathbb{C}),
\]
with the convention that the product is $0$ when $f(z)=0$, which it is, the
second factor vanishing.

\emph{Step 0: $\Omega-S$ is open, and $S$ is a null set.}  Let $z\in\Omega-S$ and
choose $\rho>0$ with $\overline{B}(z,\rho)\subseteq\Omega$.  By (H1) the set
$S\cap\overline{B}(z,\rho)$ is finite and does not contain $z$, so it has a
positive distance $\delta$ to $z$; then $B(z,\min(\rho,\delta))$ misses $S$.
Hence $\Omega-S$ is open.  Moreover $\Omega$ is a countable union of compact
subsets of $\Omega$ (for instance the closed balls $\overline{B}(q,r)\subseteq
\Omega$ with $q$ of rational coordinates and $r$ rational), so by (H1) the set
$S$ is countable and therefore Lebesgue null.

\emph{Step 1: (12) holds at every point of $\Omega-S$.}  By (H2), $f$ does not
vanish on $\Omega-S$, and by (H4) it is $C^\infty$ there; by Step 0 that set is
open.  Hence $z\mapsto\mathrm{d}f(z)$ is continuous on $\Omega-S$, so $f_z$ and
$f_{\bar z}$, being the fixed combinations
$\tfrac12(\mathrm{d}f[1]\mp i\,\mathrm{d}f[i])$ of \noderef{DZ} and
\noderef{DZbar}, are continuous there, and so are
$\overline{f}/f$ and $f/\overline{f}$.  Both sides of
\begin{equation}\tag{12}
  f_{\bar z}=c'\Bigl[\frac{\overline{f}}{f}f_z+\frac{f}{\overline{f}}
    \overline{f_z}\Bigr]
\end{equation}
are therefore continuous on $\Omega-S$.

By (H5) and \noderef{SolvesEq12} there are weak Wirtinger derivatives $a_f,b_f$
of $f$ on $\Omega$, in the sense of \noderef{IsWeakWirtingerOn}, and a null set
$Z_0$ such that
\[
  b_f(z)=c'\Bigl[\frac{\overline{f(z)}}{f(z)}a_f(z)
    +\frac{f(z)}{\overline{f(z)}}\overline{a_f(z)}\Bigr]
\]
at every $z\in\Omega-Z_0$ with $f(z)\neq0$.  We replace $a_f,b_f$ by $f_z$,
$f_{\bar z}$.  By \noderef{IsWeakWirtingerOn} there are $f_x,f_y$ with
$f,f_x,f_y\in L^1_{\mathrm{loc}}(\Omega)$ in the sense of \noderef{MemLpLocOn},
with $f_x$ and $f_y$ the distributional $\partial_0$- and
$\partial_1$-derivatives of $f$ on $\Omega$ in the sense of
\noderef{IsWeakPartialOn}, and with $a_f=\tfrac12(f_x-if_y)$,
$b_f=\tfrac12(f_x+if_y)$ at every point.  These are hypotheses (H1) and (H2) of
\noderef{WeakPartialEqFDerivAE} for the open set $\Omega$, the function $f$ and
the pair $D_0=f_x$, $D_1=f_y$; and its (H3) holds, because by (H4) and Step 0
$f$ is differentiable at every point of the open set $\Omega-S$, while $S$ is
null by Step 0.  So there is a null set $Z_1$ with $f_x=\partial_0f$ and
$f_y=\partial_1f$ on $\Omega-Z_1$, where $\partial_j=D^{[j]}$ is the classical
partial derivative of \noderef{PartialWord}, that is
$\partial_0f(z)=\mathrm{d}f(z)[1]$ and $\partial_1f(z)=\mathrm{d}f(z)[i]$.  Since
$f_z$ and $f_{\bar z}$ are by \noderef{DZ} and \noderef{DZbar} the same two
combinations $\tfrac12\bigl(\mathrm{d}f(z)[1]\mp i\,\mathrm{d}f(z)[i]\bigr)$,
we get $a_f=f_z$ and $b_f=f_{\bar z}$ on $\Omega-Z_1$.

Put $Z=Z_0\cup Z_1$, a null set.  By the last two paragraphs, (12) holds at
every $z\notin Z$ with $z\in\Omega$
and $f(z)\neq0$; on $\Omega-S$ the last condition is automatic by (H2).  Thus
(12) holds on $(\Omega-S)-Z$, which is dense in $\Omega-S$ because a null set
has empty interior, and two continuous functions on an open set agreeing on a
dense subset agree everywhere on it.  Hence (12) holds at every point of
$\Omega-S$.

\emph{Step 2: $f_{\bar z}$ is real on $\Omega-S$.}  Put $a=(\overline{f}/f)f_z$.
Then
\[
  \frac{f}{\overline{f}}\overline{f_z}
  =\overline{\Bigl(\frac{\overline{f}}{f}f_z\Bigr)}=\bar a ,
\]
so the bracket in (12) is $a+\bar a$, which is real; and $c'$ is real.  By
Step 1, $f_{\bar z}=c'(a+\bar a)$ at each point of $\Omega-S$, so
\begin{equation}\tag{1}
  \overline{f_{\bar z}}=f_{\bar z}\qquad\text{on }\Omega-S .
\end{equation}

\emph{Step 3: $w$ is continuous on $\Omega$ and $C^\infty$ on $\Omega-S$.}  On
$\Omega-S$ the function $f$ is $C^\infty$ and non-vanishing, and
$\zeta\mapsto|\zeta|^{p-2}$ is $C^\infty$ on $\mathbb{C}-\{0\}$ by the first
clause of \noderef{WirtingerModulusPowerRule}; hence $|f|^{p-2}$, and with it
$w$, is $C^\infty$ on $\Omega-S$.  For continuity on $\Omega$ it remains to
check the points of $S$.  Where $f\neq0$ we have
$|w|=|f|^{p-2}|f|=|f|^{p-1}$, and $w=0$ where $f=0$; so
$|w(z)|=|f(z)|^{p-1}$ for every $z$.  If $z_0\in S$ then $f(z_0)=0$ by (H2),
and $f$ is continuous at $z_0$ by (H3), so $|w(z)|=|f(z)|^{p-1}\to0=|w(z_0)|$
as $z\to z_0$ in $\Omega$, using $p-1>0$.  Hence $w$ is continuous on $\Omega$.

\emph{Step 4: the formula for $w_{\bar z}$ on $\Omega-S$.}  Fix $z\in\Omega-S$ and
let $F(\zeta)=|\zeta|^{p-2}$.  By \noderef{WirtingerModulusPowerRule} at the
exponent $\mu=p-2$ and the point $f(z)\neq0$, $F$ is $C^\infty$, in particular
differentiable, at $f(z)$, with
\[
  F_\zeta(f(z))=\frac{p-2}{2}\cdot\frac{|f(z)|^{p-2}}{f(z)},\qquad
  F_{\bar\zeta}(f(z))=\frac{p-2}{2}\cdot\frac{|f(z)|^{p-2}}
    {\overline{f(z)}} .
\]
Since $f$ is differentiable at $z$, \noderef{WirtingerChainRule} applies to the
pair $(F,f)$ at $z$ and gives
\[
  \partial_{\bar z}\bigl(|f|^{p-2}\bigr)
  =F_\zeta(f)f_{\bar z}+F_{\bar\zeta}(f)\overline{f_z}
  =\frac{p-2}{2}|f|^{p-2}\Bigl[\frac{f_{\bar z}}{f}
    +\frac{\overline{f_z}}{\overline{f}}\Bigr],
\]
all arguments being $z$.  Both $|f|^{p-2}$ and $f$ are differentiable at $z$,
so \noderef{WirtingerProductRule} gives
\[
  w_{\bar z}=\partial_{\bar z}\bigl(|f|^{p-2}\bigr)\,f+|f|^{p-2}f_{\bar z}
  =|f|^{p-2}\Bigl[\frac{p-2}{2}\Bigl(f_{\bar z}
    +\frac{f}{\overline{f}}\overline{f_z}\Bigr)+f_{\bar z}\Bigr]
  =|f|^{p-2}\Bigl[\frac{p}{2}f_{\bar z}
    +\frac{p-2}{2}\cdot\frac{f}{\overline{f}}\overline{f_z}\Bigr].
\]
Multiplying the bracket by $|f|^2/|f|^2$ and using
$|f|^2=f\overline{f}$, so that $(f/\overline{f})|f|^2=f^2$, we obtain
\begin{equation}\tag{2}
  w_{\bar z}=|f|^{p-4}X,\qquad
  X=\frac p2|f|^2f_{\bar z}+\frac{p-2}{2}f^2\overline{f_z},
\end{equation}
at every point of $\Omega-S$.

\emph{Step 5: $\operatorname{Re}w_{\bar z}=0$ on $\Omega-S$.}  Fix $z\in\Omega-S$.
Using (1) and the reality of $|f|^2$,
\[
  \bar X=\frac p2|f|^2f_{\bar z}+\frac{p-2}{2}\overline{f}^{\,2}f_z ,
\]
so
\[
  X+\bar X=p|f|^2f_{\bar z}
    +\frac{p-2}{2}\bigl(f^2\overline{f_z}+\overline{f}^{\,2}f_z\bigr).
\]
By Step 1 and $c'=(2-p)/(2p)$,
\[
  p|f|^2f_{\bar z}
  =p\,f\overline{f}\cdot\frac{2-p}{2p}\Bigl[\frac{\overline{f}}{f}f_z
    +\frac{f}{\overline{f}}\overline{f_z}\Bigr]
  =\frac{2-p}{2}\bigl(\overline{f}^{\,2}f_z+f^2\overline{f_z}\bigr),
\]
the two cancellations being legitimate because $f(z)\neq0$.  Hence
\[
  X+\bar X=\frac{2-p}{2}\bigl(\overline{f}^{\,2}f_z+f^2\overline{f_z}\bigr)
    +\frac{p-2}{2}\bigl(f^2\overline{f_z}+\overline{f}^{\,2}f_z\bigr)=0,
\]
the two coefficients being exact negatives.  Since $X+\bar X=2\operatorname{Re}X$
we get $\operatorname{Re}X=0$, and by (2), $|f|^{p-4}$ being real,
\begin{equation}\tag{3}
  \operatorname{Re}w_{\bar z}=0\qquad\text{on }\Omega-S .
\end{equation}

\emph{Step 6: the gradient field.}  Define
\[
  g(z)=2\overline{f(z)}\qquad(z\in\mathbb{C}).
\]
On $\Omega-S$ we have, by (H8) and \noderef{DZ},
$\tfrac12\bigl(\partial_x\varphi_v-i\,\partial_y\varphi_v\bigr)=f$; the function
$v$ is real-valued, so $\partial_x\varphi_v$ and $\partial_y\varphi_v$ are
real, and comparing real and imaginary parts,
\begin{equation}\tag{4}
  \partial_x\varphi_v=2\operatorname{Re}f=\operatorname{Re}g,\qquad
  \partial_y\varphi_v=-2\operatorname{Im}f=\operatorname{Im}g
  \qquad\text{on }\Omega-S .
\end{equation}
Next, at every $z\in\mathbb{C}$,
\begin{equation}\tag{5}
  |g|^{p-2}\operatorname{Re}g=2^{p-1}\operatorname{Re}w,\qquad
  |g|^{p-2}\operatorname{Im}g=-2^{p-1}\operatorname{Im}w .
\end{equation}
Indeed, if $f(z)\neq0$ then $|g|=2|f|\neq0$ and
$|g|^{p-2}\operatorname{Re}g=2^{p-2}|f|^{p-2}\cdot2\operatorname{Re}f
=2^{p-1}\operatorname{Re}\bigl(|f|^{p-2}f\bigr)$, and similarly for the
imaginary part; and if $f(z)=0$ then $g(z)=0$, so both sides of each identity
vanish, the left ones because the factors $\operatorname{Re}g$ and
$\operatorname{Im}g$ do.

\emph{Step 7: the weak gradient of $v$ on all of $\Omega$.}  We apply
\noderef{WeakPartialAcrossPoint} on $\Omega$ with the exceptional set $S$,
which satisfies its local finiteness hypothesis by (H1), and with growth
exponent $\beta=0$.  Take first $j=0$, the function $\varphi_v$ and the
candidate derivative $\operatorname{Re}g$.  Its hypotheses hold:
$\varphi_v$ is $C^1$ on $\Omega-S$ by (H7); its classical $x$-partial there is
$\operatorname{Re}g$ by (4); $\varphi_v$ is continuous on $\Omega$ by (H6) and
$\operatorname{Re}g$ is continuous on $\Omega$ by (H3), so both are bounded on
every compact subset of $\Omega$ and hence lie in
$L^1_{\mathrm{loc}}(\Omega)$; and for $z_0\in S$, continuity of $\varphi_v$ at
$z_0$ gives $M,\rho>0$ with $|\varphi_v(z)|\le M=M|z-z_0|^{-0}$ for
$0<|z-z_0|<\rho$.  The conclusion is that $\operatorname{Re}g$ is the
distributional $x$-partial of $\varphi_v$ on $\Omega$ in the sense of
\noderef{IsWeakPartialOn}.  Taking $j=1$ with the candidate
$\operatorname{Im}g$ gives in the same way that $\operatorname{Im}g$ is the
distributional $y$-partial of $\varphi_v$ on $\Omega$.

\emph{Step 8: the weak form of the $p$-harmonic equation.}  Let
$W_1=\operatorname{Re}w$ and $W_2=\operatorname{Im}w$, regarded as
complex-valued functions.  By Step 3 they are continuous on $\Omega$ and
$C^\infty$ on $\Omega-S$.  Apply \noderef{WeakPartialAcrossPoint} on $\Omega$
with exceptional set $S$, $\beta=0$, first to $j=0$, the function $W_1$ and its
classical $x$-partial $\partial_xW_1$ on $\Omega-S$, and then to $j=1$, the
function $W_2$ and its classical $y$-partial $\partial_yW_2$.  The local
integrability of $W_1,W_2$ follows from their continuity on $\Omega$; the local
integrability of $\partial_xW_1$ and $\partial_yW_2$ follows from (H10):
taking real or imaginary parts is a continuous $\mathbb{R}$-linear operation and
so commutes with differentiation, so on $\Omega-S$ we have
$\partial_xW_1=\operatorname{Re}(\partial_0w)$ and
$\partial_yW_2=\operatorname{Im}(\partial_1w)$, dominated in modulus by
$|\partial_0w|$ and $|\partial_1w|$ respectively, and $S$ is null by Step 0, so
the domination off a null set suffices for integrability on compact subsets of
$\Omega$; and the growth hypothesis
holds with $\beta=0$ because $W_1,W_2$ are continuous, hence bounded near each
point of $S$.  We obtain
\begin{equation}\tag{6}
  \int W_1\,\partial_x\phi=-\int(\partial_xW_1)\phi,\qquad
  \int W_2\,\partial_y\phi=-\int(\partial_yW_2)\phi
\end{equation}
for every test function $\phi$ for $\Omega$ in the sense of \noderef{IsTestFn}.

On $\Omega-S$ we compute, from \noderef{DZbar},
\[
  \operatorname{Re}w_{\bar z}
  =\operatorname{Re}\Bigl(\tfrac12(\partial_xw+i\,\partial_yw)\Bigr)
  =\tfrac12\bigl(\partial_x\operatorname{Re}w-\partial_y\operatorname{Im}w\bigr)
  =\tfrac12\bigl(\partial_xW_1-\partial_yW_2\bigr),
\]
so by (3),
\begin{equation}\tag{7}
  \partial_xW_1-\partial_yW_2=0\qquad\text{on }\Omega-S .
\end{equation}

Now let $\phi$ be a real-valued test function for $\Omega$; regarded as
complex-valued it is a test function in the sense of \noderef{IsTestFn}.  Using
(5), then (6), then (7) together with the fact that $S$ is null by Step 0 and
that $\phi$ vanishes off a compact subset of $\Omega$,
\[
  \int|g|^{p-2}\bigl(\operatorname{Re}g\cdot\partial_x\phi
    +\operatorname{Im}g\cdot\partial_y\phi\bigr)
  =2^{p-1}\int\bigl(W_1\partial_x\phi-W_2\partial_y\phi\bigr)
  =-2^{p-1}\int\bigl(\partial_xW_1-\partial_yW_2\bigr)\phi=0 .
\]
Finally, since $\partial_x\phi$ and $\partial_y\phi$ are real, at every $z$
\[
  \operatorname{Re}\Bigl(g(z)\,
    \overline{\partial_x\phi(z)+i\,\partial_y\phi(z)}\Bigr)
  =\operatorname{Re}g(z)\,\partial_x\phi(z)
    +\operatorname{Im}g(z)\,\partial_y\phi(z).
\]
Hence
\[
  \int|g|^{p-2}\operatorname{Re}
    \Bigl(g\,\overline{\partial_x\phi+i\,\partial_y\phi}\Bigr)=0
\]
for every real test function $\phi$ for $\Omega$, which is the equation clause
of \noderef{PHarmonicOn}.

\emph{Step 9: conclusion.}  Take $g=2\overline{f}$ as the witness required by
\noderef{PHarmonicOn}.  Then $v\in L^p_{\mathrm{loc}}(\Omega)$ by (H9); and
$g\in L^p_{\mathrm{loc}}(\Omega)$ because $|g|=2|f|$ and
$f\in L^p_{\mathrm{loc}}(\Omega)$ by (H9).  Step 7 supplies the two weak
partial derivatives of $v$, namely $\operatorname{Re}g$ in the $x$ direction
and $\operatorname{Im}g$ in the $y$ direction, and Step 8 supplies the
equation.  Hence $v$ is $p$-harmonic on $\Omega$ in the sense of
\noderef{PHarmonicOn}.
\end{proof}
